\question[4]
Wir vergleichen zwei Vorgehensmodelle.
\begin{parts}
\part Beschreibe den Unterschied zwischen dem Wasserfallmodell und dem agilen Vorgehen.
\part Eine Firma lässt eine App entwickeln, weiß aber noch nicht genau, welche Funktionen die Kunden brauchen. Welches Vorgehen empfiehlst du? Begründe.
\end{parts}
\begin{solutionbox}{4cm}
a) Beim Wasserfallmodell laufen die Phasen einmal nacheinander ab, die Anforderungen stehen am Anfang fest. Beim agilen Vorgehen wird in kurzen Runden gearbeitet, nach jeder Runde gibt es ein lauffähiges Zwischenergebnis. \\
b) agiles Vorgehen, denn die Anforderungen sind noch unklar. Die Firma sieht früh Zwischenergebnisse und kann die Anforderungen anpassen.
\end{solutionbox}
